\section{Marked stable realization}\label{st:setup}

We now suppose, for a contradiction, that the group constructed above is
the fundamental group of a closed aspherical topological four-manifold
$M$.  The preceding construction gives a finite oriented Poincar\'e
four-complex $Q=BG$ and a closed oriented manifold $N$.  Write $q$ for
the map of Proposition~\ref{rf:development}:
\[
 q\colon N\longrightarrow Q
\]
which is an isomorphism on fundamental groups and has degree one.  The
only reduced homology of the universal cover of $N$ is its second
homology.  With $\Lambda=\Z G$, this is the free module
\begin{equation}\label{st:module}
 A=\pi_2(N)\cong\Lambda^{2b},
 \qquad \lambda_N\cong\bigoplus_{a=1}^{b}
 \begin{pmatrix}0&1\\1&0\end{pmatrix}.
\end{equation}
The displayed basis consists of the prescribed chamber sphere pairs.
All identifications include whiskers and orientations.  Moreover, $G$
is torsion-free, is an integral $\mathrm{PD}_4$ group, and satisfies the
$K$- and $L$-theoretic Farrell--Jones conjectures by the CAT(0)
construction.  These are conclusions of the construction, not
additional assumptions on a putative realization.

Both $Q$ and $M$ are $K(G,1)$ spaces.  Their equivalence identifies the
orientation character with the action on the top dualizing module.
Consequently $M$ is orientable.  Choose its orientation and its
fundamental-group marking so that its classifying map
$c_M\colon M\to Q$ has degree one, and put
\[
 L=M\mathbin{\#}\,b(S^2\times S^2).
\]
Collapsing the simply connected summands gives a degree-one classifying
map $c_L\colon L\to Q$.  Identify $\pi_2(L)$ with $A$ by sending each
specified chamber pair in \eqref{st:module} to a specified standard pair
in $L$.  This is an isometry of equivariant forms.  We must retain this
particular isometry throughout the argument.

\subsection{An integral marked homotopy equivalence}

We first realize the prescribed module isometry by a homotopy
equivalence. The comparison takes place in a common second Postnikov
stage: group duality splits that stage, and integral cup-product tests
identify the two fundamental classes there. These are the inputs to
the lifting argument below.

\begin{lemma}\label{st:second-stage}
The second Postnikov stages of $N$ and $L$, with the preceding markings,
identify over $Q$ with the split stage
\[
 B=EG\times_G K(A,2).
\]
There are maps $u_N\colon N\to B$ and $u_L\colon L\to B$ inducing the
specified identifications of fundamental and second homotopy groups.
If $\Gamma(A)=H_4(K(A,2);\Z)$, then
\begin{equation}\label{st:homology-extension}
 0\longrightarrow\Gamma(A)_G\longrightarrow H_4(B;\Z)
 \longrightarrow H_4(Q;\Z)\longrightarrow0
\end{equation}
is a split exact sequence.
\end{lemma}

\begin{proof}
Since $A$ is finite free over $\Lambda$, integral group duality gives
\begin{equation}\label{st:cohomology-vanishing}
 H^2(G;A)=H^3(G;A)=0.
\end{equation}
The second Postnikov invariant lies in the latter group, so it vanishes.
The former removes the degree-two ambiguity in making the identification
with the split stage for the prescribed module marking.  We use a
functorial Eilenberg--Mac Lane model for $K(A,2)$, so the given action on
$A$ defines the displayed homotopy quotient and its section.

Separately, freeness gives $H_i(G;A)=0$ for $i>0$.  In total degree four,
the homology spectral sequence of
\[
 K(A,2)\longrightarrow B\longrightarrow Q
\]
therefore has only $H_4(G;\Z)$ and $H_0(G;\Gamma(A))$ as possibly
nonzero terms: the intervening term $H_2(G;A)$ is zero, and the fiber has
no homology in degrees one and three.  The possible incoming
differentials to $H_0(G;\Gamma(A))$ have sources $H_3(G;A)=0$ and
$H_5(G;\Z)=0$.  The section preserves the base term.  This proves
\eqref{st:homology-extension}, including its integral, rather than merely
rational, exactness.
\end{proof}

The next elementary check is useful because the middle classes are
specified over the group ring.  Ordinary intersection forms would lose
the relative translates needed here.

\begin{lemma}\label{st:quadratic-detection}
With the identifications above,
\begin{equation}\label{st:fundamental-equality}
 (u_N)_*[N]=(u_L)_*[L]\quad\text{in }H_4(B;\Z).
\end{equation}
\end{lemma}

\begin{proof}
Write $e_1,\ldots,e_{2b}$ for the chosen $\Lambda$-basis of $A$ and
$\phi_i\colon A\to\Lambda$ for the coordinate functionals.  The
low-degree cohomology sequence of the Postnikov fibration, with regular
coefficients, contains
\[
 0\longrightarrow H^2(G;\Lambda)
 \longrightarrow H^2(B;\Lambda)
 \longrightarrow\operatorname{Hom}_{\Lambda}(A,\Lambda)
 \longrightarrow H^3(G;\Lambda).
\]
Both outside groups vanish.  Thus each $\phi_i$ is the restriction of
a unique class $x_i\in H^2(B;\Lambda)$.

Their cup products may be evaluated without multiplying away either
coefficient factor:
\begin{equation}\label{st:cup-tests}
 z\longmapsto
 \langle x_i\smile x_j,z\rangle
 \in(\Lambda\otimes_{\Z}\Lambda)_G,
 \qquad z\in H_4(B;\Z).
\end{equation}
The tensor product has the diagonal left action.  Its coinvariants
identify, as an abelian group, with $\Lambda$ by
$g\otimes h\mapsto g^{-1}h$.  For a closed oriented manifold with
the indicated middle module, evaluations in \eqref{st:cup-tests} are
the equivariant cup form on the coordinate functionals.  By Poincar\'e
duality this is the dual form of its equivariant intersection pairing.
The prescribed isometry therefore makes these evaluations equal on
$(u_N)_*[N]$ and $(u_L)_*[L]$.

For completeness, these tests detect the fiber term integrally.  The
underlying abelian group of $A$ is free on the elements $g e_i$.
The group $\Gamma(A)$ has a basis consisting of one square generator
for each such element and one cross generator for each unordered pair
of distinct such elements.  The action of $G$ permutes this basis.
Hence its coinvariants are free abelian on the corresponding orbits.
A nonidentity element stabilizing an unordered pair would exchange its
two members and have order two; torsion-freeness rules this out.

The square orbit represented by $e_i$ evaluates to $1$ in the $(i,i)$
test.  For $i<j$, the cross orbit represented by
$\{e_i,t e_j\}$ evaluates to $t$ in the $(i,j)$ test.  Finally, for
$t\ne1$, the orbit represented by $\{e_i,t e_i\}$ evaluates to
$t+t^{-1}$ in the $(i,i)$ test.  These last orbits are indexed by $t$
modulo inversion.  Since $G$ has no elements of order two, their supports
are disjoint two-element subsets of $G\setminus\{1\}$.  The displayed
tests are therefore jointly injective on $\Gamma(A)_G$.  In particular,
no integral square or hidden two-torsion is discarded.

The two fundamental classes have the same image $[Q]$ in $H_4(Q;\Z)$.
By Lemma~\ref{st:second-stage}, their difference lies in
$\Gamma(A)_G$.  All its cup tests vanish, so the difference is zero.
\end{proof}

\begin{proposition}\label{st:marked-equivalence}
There is an orientation-preserving homotopy equivalence
\[
 f\colon N\longrightarrow L
\]
inducing the specified isomorphism on $\pi_1$ and the prescribed
isometry $A\to\pi_2(L)$.
\end{proposition}

\begin{proof}
Replace $u_L\colon L\to B$ by a fibration.  Its fiber is
two-connected, with third homotopy group $\pi_3(L)$ carrying the
usual $G$-action.  The first obstruction to lifting $u_N$ is consequently
the pullback of a universal class
\[
 \kappa\in H^4(B;\pi_3(L)).
\]
The pullback $u_L^*\kappa$ is zero, since $u_L$ lifts to its own
fibration replacement.  Naturality of cap product and
Lemma~\ref{st:quadratic-detection} give
\begin{align*}
 (u_N)_*\bigl((u_N^*\kappa)\mathbin{\cap}[N]\bigr)
 &=\kappa\mathbin{\cap}(u_N)_*[N]\\
 &=\kappa\mathbin{\cap}(u_L)_*[L]=0
 \quad\text{in }H_0(B;\pi_3(L)).
\end{align*}
The map $u_N$ induces an isomorphism on fundamental groups, and hence
on these degree-zero coefficient groups.  Poincar\'e duality identifies
\[
 H^4(N;\pi_3(L))\xrightarrow{\ \cap[N]\ }
 H_0(N;\pi_3(L))\cong\pi_3(L)_G.
\]
Thus the lifting obstruction vanishes.

The manifold $N$ has a finite four-dimensional CW structure from the
smooth chamber construction and the panel collars.  There are no
higher obstructions on this source.  No triangulation or smooth
structure on the hypothetical $M$ is required.  The resulting lift
$f$ has the stipulated maps on $\pi_1$ and $\pi_2$.  The universal
covers of $N$ and $L$ have reduced homology only in degree two, where
the lift is an isomorphism. These manifolds have CW homotopy type
\cite[Corollary~1]{MilnorCW}. The homological Whitehead theorem on their
simply connected universal covers, followed by Whitehead's theorem,
shows that $f$ is a homotopy equivalence
\cite[Corollary~4.33 and Theorem~4.5]{Hatcher}.
Finally $c_Lf\simeq q$ and both classifying maps have degree one, so
$f_*[N]=[L]$.
\end{proof}

\subsection{From the marked equivalence to a marked stable product}

The marked homotopy equivalence is now available. The next two steps
pass through an $s$-cobordism and then a stabilized product, retaining
the same identification of every original sphere class. The product
trace also specifies the new classes introduced by stabilization.

\begin{lemma}\label{st:marked-bordism}
There is an $s$-cobordism $(W;N,L)$ whose two end inclusions identify
$\pi_1$ and $\pi_2$ by the marking of
Proposition~\ref{st:marked-equivalence}.
\end{lemma}

\begin{proof}
The universal intersection forms are even, so $N$ and $L$ are almost
spin.  Indeed, evaluation of $w_2$ on every spherical class is its
self-intersection modulo two, and on the simply connected cover these
evaluations detect $w_2$.

We apply the proof of Kasprowski--Land's Theorem~1.1, specifically
Lemma~2.3, Lemma~3.3 and Theorem~3.1 of~\cite{KL}.  All hypotheses hold:
$Q$ is an aspherical Poincar\'e four-complex; $G$ satisfies both
Farrell--Jones conjectures; the two manifolds have the same orientation
character and degree-one classifying maps; and they are almost spin.
Their normal-invariant calculation, using assembly injectivity and the
Atiyah--Hirzebruch spectral sequence, identifies the normal invariant of
this particular $f$ with a class in
\[
 \ker\bigl(H_2(L;\Z/2)\longrightarrow H_2(Q;\Z/2)\bigr).
\]
The Serre spectral sequence for the universal-cover fibration shows
that this kernel is represented by spherical classes. Since $L$ is
almost spin, their $w_2$ evaluations vanish. Lemma~3.3 of~\cite{KL}
therefore supplies a target self-equivalence $\phi$ with the same
normal invariant as $f$, of the form
\[
 L\longrightarrow L\vee S^4
 \xrightarrow{\operatorname{id}\vee\eta^2}L\vee S^2
 \xrightarrow{\operatorname{id}+a}L.
\]
Here $a:S^2\to L$ represents the spherical class. This pinch acts
identically on $\pi_1$ and $\pi_2$, because its extra part factors through
the four-sphere. Thus $\phi^{-1}\circ f$ has trivial normal invariant
and the same marking as $f$. Surgery of its normal bordism over $L$,
as in Theorem~3.1 of~\cite{KL}, gives the required $s$-cobordism.
Neither this conclusion nor
the cited theorem requires $G$ to be a good group.  We have not asserted
that $W$ itself is a product.
\end{proof}

We use the stable product theorem for this particular bordism.  Quinn's
Theorem~1.1 in~\cite{QuinnStable}, in the form stated in
\cite[Remark~3.7]{KPR}, says that a five-dimensional topological
$s$-cobordism becomes a product relative to one end after finitely many
connected sums along arcs with $(S^2\times S^2)\times[0,1]$.
There is no restriction on its fundamental group.  The relative product
statement is important: an arbitrary stable homeomorphism of the ends
would not, by itself, retain the specified sphere classes.

\begin{proposition}\label{st:stable-product}
For some integer $k\ge0$ there is a homeomorphism
\[
 h\colon N\mathbin{\#}\,k(S^2\times S^2)
 \longrightarrow
 L\mathbin{\#}\,k(S^2\times S^2)
\]
which induces the prescribed original marking and matches every added
standard pair with its specified added standard pair.  In particular,
on the left, all the prescribed original hyperbolic pairs and all new
standard pairs have simultaneous geometric representatives with
pairwise disjoint punctured $S^2\times S^2$ neighborhoods.
\end{proposition}

\begin{proof}
Apply the stable product theorem to $W$ from
Lemma~\ref{st:marked-bordism}.  Denote its stabilized ends by $N^*$ and
$L^*$ and the stabilized bordism by $W^*$.  We verify its marking,
including the new summands.

The operation removes disjoint normal $D^4\times[0,1]$ neighborhoods
of arcs between the ends and glues
\[
 \bigl((S^2\times S^2)\setminus\operatorname{int}D^4\bigr)
 \times[0,1]
\]
along the resulting $S^3\times[0,1]$ interfaces.  Removing these
codimension-four arcs changes neither $\pi_1$ nor second homology of
the universal cover.  Mayer--Vietoris across the lifted interfaces
therefore gives a specified decomposition
\begin{equation}\label{st:bordism-splitting}
 \pi_2(W^*)\cong\pi_2(W)\oplus\Lambda^{2k}.
\end{equation}
The second summand is represented by the standard spheres in the added
product pieces.  Choose their whiskers along the stabilization arcs.
The two end inclusions $j_N\colon N^*\to W^*$ and
$j_L\colon L^*\to W^*$ carry each new pair to this same pair, while
their old summands map to $\pi_2(W)$ with exactly the original bordism
identification.

A product structure $F\colon N^*\times[0,1]\to W^*$ which is the
identity on the incoming end restricts at the other end to a
homeomorphism $h\colon N^*\to L^*$.  Use the product trace for the
outgoing basepoint path.  Any discrepancy from the original bordism
path can be corrected by pushing the outgoing basepoint around a loop
in $L^*$, since $j_L$ is an isomorphism on fundamental groups.
The resulting based homotopy given by $F$ implies
\[
 (j_L)_*h_*=(j_N)_*.
\]
Both inclusions are isomorphisms on $\pi_1$ and $\pi_2$.  Thus
$h_*=(j_L)_*^{-1}(j_N)_*$, with the same base paths.  In the
decomposition \eqref{st:bordism-splitting} it
preserves the original marking and identifies the new standard basis
element by element.  There is no arbitrary automorphism mixing the
old and new summands.

All $b+k$ standard pairs on
$L^*=M\mathbin{\#}(b+k)(S^2\times S^2)$ have disjoint punctured
$S^2\times S^2$ neighborhoods.  Give the two spheres of each pair
the common whisker to their intersection used for its module marking,
and pull the whole based neighborhood back by $h$.  A pair has
its one prescribed transverse intersection; different pairs and their
neighborhoods are disjoint.  The marking just proved identifies their
based homotopy classes with the original chamber basis and the added
standard basis, respectively.
\end{proof}

\subsection{Periodically lifted geometric data}

Let $C$ be the right-angled Coxeter group used in the reflection
construction, and let $\Delta\le C$ be the finite-index torsion-free
normal subgroup used to obtain $N$ from its ordinary reflection
development.  Place the finitely many stabilization balls on the
$N$-side inside chamber interiors.  This may be done by moving their
endpoints along paths before performing the arc connected sums.
Track the new whiskers along those paths.  There is no requirement that
different chambers receive equal numbers of stabilizations.

Lift the resulting $N^*$ to the ordinary reflection development and
call the lifted manifold $Y$.  Its chambers, indexed by $v\in C$, are
stabilized copies $D_v$ of $D$.  The numbers of stabilizations are
$\Delta$-periodic.  Gallery distance is the word metric on the chamber
indices for the standard Coxeter generators.

\begin{proposition}\label{st:periodic-data}
The manifold $Y$ has two $\Delta$-periodic systems of sphere pairs,
with the same chamber labels:
\begin{enumerate}
\item framed immersed hyperbolic pairs in the interiors of the
chambers $D_v$, including the added standard pairs, with their
equivariant products already correct over the individual chamber group;
\item geometrically standard pairs with pairwise disjoint punctured
$S^2\times S^2$ neighborhoods, based-homotopic sphere by sphere to
the correspondingly labelled pairs in the first system.
\end{enumerate}
There is an integer $d_0$ such that the drawing of each pair, its
standard neighborhood, and the homotopies relating the two systems
are contained in the chambers at gallery distance at most $d_0$ from
that pair's label.  Basings in this statement are local basings
transported with the chamber; no bound is asserted for whiskers drawn
from one global base point to all chambers.
\end{proposition}

\begin{proof}
The original immersed representatives stay in their chamber interiors,
and the stabilization spheres can be put in their interior connected
summands.  Stabilization adds orthogonal standard planes and changes
neither the chamber fundamental group nor the old pairings.

Proposition~\ref{st:stable-product} supplies a finite collection of
geometric pairs on the compact quotient $N^*$, representing the
specified based classes.  Choose a based homotopy for each sphere
between its two representatives.  Lift the neighborhoods, sphere
drawings, and homotopies to $Y$, using the prescribed marking to choose
the lift for each chamber label.  Covering maps preserve $\pi_2$ and
lift these based homotopies, so they end at precisely the labelled
geometric representatives.  Disjoint embedded neighborhoods lift to
disjoint embedded neighborhoods.

There are only finitely many labelled data on the quotient.  Each chosen
lift of a drawing, closed standard neighborhood, local basing, or
homotopy has compact image.  Local finiteness of the development implies
that it meets only finitely many chambers.  Take $d_0$ to be the maximum
gallery distance from its assigned label over this finite list of
chosen lifts.  All remaining data are translates by $\Delta$, whose
action preserves gallery distance.  The same $d_0$ therefore works
everywhere.
\end{proof}

The last proposition is the stable information needed below.  It gives
actual geometric summands with their original chamber labels and
compactly supported homotopies to the chamber systems.  It makes no
claim of controlled isotopy, and no claim that the unstabilized chamber
obstruction has disappeared.
